\section{Significance of Research}
% =============================================================================
% SIGNIFICANCE OF RESEARCH
% =============================================================================
% REQUIREMENT: Demonstrate importance by explaining why topic is relevant,
%              how findings will be relevant, why timely to research now
%
% STATUS: Draft written - captures main points but could be stronger
%
% CURRENT TEXT COVERS:
%   - Ongoing debate (relevance) ✓
%   - ABM widely used, AI approach novel (contribution) ✓
%   - RL vs LLM comparison is active research area (timeliness) ✓
%
% REMAINING TODOS:
%   - Fix typo: "comparitive" → "comparative"
%   - Strengthen "could potentially provide" → more confident claim
%   - Be more specific about YOUR contribution vs general topic importance
%   - Consider adding: "First direct computational comparison of neorealism
%     vs constructivism" as explicit novelty claim
%   - Consider citing: Cederman 2003 (constructivists need formal models),
%     Unver 2018 (computational IR gap), Larooij 2025 (generative ABM value)
%
% STRONGER FRAMING OPTIONS:
%   1. "This research provides the first computational framework for directly
%      comparing neorealist and constructivist predictions in the same simulation"
%   2. "By operationalizing theoretical assumptions as different agent
%      architectures, this work enables empirical testing of previously
%      untestable theoretical claims"
%   3. "The RL-LLM comparison addresses calls for formal modeling tools in
%      constructivist IR (Cederman & Daase 2003) while leveraging recent
%      advances in generative agents"
% =============================================================================

With the debate between Neorealism and Constructivism still actively ongoing, my research could potentially provide additional insights. Though, ABM models are widely used in IR studies, the Artificial Intelligence approach to designing the agent architecture seems quite novel. Especially, comparitive studies between RL and LLM behaviour in social studies is an active field of research.
